\section{Related Work}\label{sec:rw}

% Skeleton: five sub-areas mapped to the foundations the approach
% builds on. Citations are pre-loaded; prose to be expanded.

\paragraph{Trace-link recovery: IR foundations.}
Automated traceability link recovery emerged as a sub-field around the requirements-traceability problem statement of Gotel and Finkelstein~\cite{gotel1994analysis} and the link-type reference models of Ramesh and Jarke~\cite{ramesh2001toward}; we use the consolidated definitions of Cleland-Huang, Gotel, and Zisman~\cite{clelandhuang2012handbook} and the roadmap of Spanoudakis and Zisman~\cite{spanoudakis2005software} throughout.
The dominant recovery family for two decades was information retrieval over textual artefacts: the vector-space and probabilistic models of Antoniol et al.~\cite{antoniol2002recovering}, latent-semantic indexing as introduced to traceability by Marcus and Maletic~\cite{marcus2003recovering}, the artefact-management variants of De Lucia et al.~\cite{delucia2007recovering,delucia2009assessing}, and the candidate-ranking framework of Hayes, Dekhtyar, and Sundaram~\cite{hayes2006advancing}.
Borg, Runeson, and Ard{\"o}~\cite{borg2014recovering} systematically map the first decade of IR-based recovery and document its evaluation heterogeneity; M{\"a}der and Egyed~\cite{maeder2015developers} provide the strongest empirical case that the recovered links pay back in maintenance.

\paragraph{Bridging the vocabulary gap.}
The recurring obstacle across IR-based recovery is the vocabulary mismatch first articulated by Furnas et al.~\cite{furnas1987vocabulary}: the same referent surfaces under many forms.
Two responses dominate.
The concept- and feature-location lines inject human-curated vocabulary into the retrieval step~\cite{rajlich2002role,dit2013feature}; noun-anchored indexing~\cite{capobianco2013role} and semantic enhancement~\cite{mahmoud2015role} are SE-traceability-internal moves in the same spirit.
The requirements-engineering side instead mines vocabulary from the text itself: glossary-term extraction~\cite{arora2017automated,gemkow2018automatic}, abbreviation expansion and identifier normalisation in code~\cite{lawrie2011expanding}, mining synonyms from comment-code mappings~\cite{howard2013automatically}, and WordNet-like identifier networks~\cite{falleri2010automatic}.
\approach extends this line by treating an alias table and an ambiguity map as a first-class artefact built before any link is scored, rather than as a query- or index-time preprocessing step.

\paragraph{Entity linking, coreference, and anaphora.}
Outside software engineering, the entity-linking community has developed the surface-mention view of linking that \approach inherits: a link from a textual span to a knowledge-base entity is asserted only when the span surfaces as a recognised reference~\cite{mihalcea2007wikify,bunescu2006using,cucerzan2007large,hachey2013evaluating}.
Abbreviation discovery~\cite{schwartz2003simple} and pronoun resolution~\cite{hobbs1978resolving,lappin1994algorithm,grosz1995centering,mitkov1998robust,soon2001machine} provide the classical machinery for the alias and pronominal cases.
The requirements-engineering community has, more recently, brought anaphora resolution to bear on SE artefacts directly~\cite{yang2011analysing,ezzini2022automated}, but, to our knowledge, no prior trace-link recovery approach uses pronominal references as link evidence.

\paragraph{Architecture-to-code traceability.}
Architecture documents differ from requirements documents in referential structure: components are named entities, prose mixes canonical names with project-specific aliases, and the documentation-code consistency problem is itself a long-running sub-field~\cite{murphy2001software}.
ArchTrace~\cite{murta2008continuous} and the LISA toolkit~\cite{buchgeher2011automatic} are early architecture-side trace-recovery tools; Mirakhorli and Cleland-Huang~\cite{mirakhorli2016detecting} link architectural tactics to code; Javed and Zdun~\cite{javed2014systematic} summarise the sub-field; Tian, Liang, and Babar~\cite{tian2021relationships} document the practitioner view.
The ARDoCo line~\cite{keim_trace_2021,keim_recovering_2024} establishes the benchmark and the SAD-SAM-code pipeline that this work compares against, and LiSSA~\cite{fuchs_lissa_2025} demonstrates that retrieval-augmented LLMs can replace several IR-era stages.

\paragraph{Position of \approach.}
\approach is the first architecture-to-code trace-link recovery approach that (i)~separates the runtime knowledge layer from the linker as a first-class architectural decision, (ii)~uses pronominal references as link evidence under a structural antecedent constraint, and (iii)~matches \ac{LLM} validation pass count to the epistemic difficulty of each sub-problem.
The two prior baselines we compare against are SWATTR~\cite{swattr}, which represents the lexical-pattern line, and LiSSA~\cite{lissa}, which represents the single-pass-LLM line.
